\documentclass{article}
\newcommand{\benchname}{BenchName}
\usepackage{neurips_2026}

\usepackage[utf8]{inputenc}
\usepackage[T1]{fontenc}
\usepackage{hyperref}
\usepackage{url}
\usepackage{booktabs}
\usepackage{amsfonts}
\usepackage{amsmath}
\usepackage{amssymb}
\usepackage{amsthm}
\usepackage{nicefrac}
\usepackage{microtype}
\usepackage{xcolor}
\usepackage[table]{xcolor}
\usepackage{colortbl}
\usepackage{multirow}
\usepackage{arydshln}
\usepackage{makecell}
\usepackage{graphicx}
\usepackage{subfigure}
\usepackage[hidelinks]{hyperref}
\usepackage{enumitem}
\usepackage[flushleft]{threeparttable}
\usepackage{subcaption}
\usepackage[table]{xcolor}
\usepackage{booktabs}
\usepackage{graphicx}
\usepackage{pifont}
\usepackage{microtype}      % microtypography
\usepackage{xcolor}         % colors
\usepackage[table]{xcolor}
\usepackage{colortbl}
\usepackage{booktabs}
\usepackage[table]{xcolor}
\usepackage{multirow}
\usepackage{arydshln}
% \newcommand{\TODO}[1]{\textcolor{red}{\textbf{[TODO: #1]}}}
\usepackage{makecell}
%\usepackage[frozencache,cachedir=.]{minted}
% Recommended, but optional, packages for figures and better typesetting:
\usepackage{microtype}
\usepackage{graphicx}
\usepackage{subfigure}
\usepackage{booktabs} % for professional tablesa
\usepackage[hidelinks]{hyperref} 
\usepackage{arydshln}
\usepackage{booktabs}
\usepackage{xspace}
\usepackage{enumitem}
\usepackage[flushleft]{threeparttable}
\usepackage{makecell}
\usepackage{multirow}
\usepackage{xcolor, colortbl}
\usepackage{subcaption}
\usepackage{tikz}
\usepackage{wrapfig}
\usepackage{textcomp}  %
\usepackage{scalerel}  %
\usepackage{tipa}
\usepackage{pifont}
\usepackage[subtle]{savetrees}
\usepackage[most]{tcolorbox}
\usepackage{listings}
\usepackage{caption}
\usepackage[normalem]{ulem}
\usepackage{tcolorbox}
\usepackage{booktabs}
\usepackage{multirow}
\usepackage{makecell}
\usepackage{array}
\usepackage{graphicx}
\usepackage[table]{xcolor}
\usepackage{colortbl}
\usepackage{bbm}
\usepackage{algorithmic}
\usepackage[table]{xcolor}
\usepackage{tabularray}
\usepackage{float}
\definecolor{HeaderGray}{gray}{0.95}
\definecolor{GroupGreen}{RGB}{220,240,220}
\definecolor{GroupBlue}{RGB}{220,230,245}
\definecolor{HeaderGray}{gray}{0.95}
\definecolor{GroupGreen}{RGB}{220,240,220}
\definecolor{GroupBlue}{RGB}{220,230,245}
\definecolor{OursRowBg}{RGB}{223, 233, 251}

\definecolor{cvprblue}{rgb}{0.21,0.49,0.8}
\definecolor{HeaderGray}{RGB}{235,238,242}
\definecolor{OursRowBg}{RGB}{242,247,255}
\definecolor{GroupBlue}{RGB}{55,105,180}
\definecolor{myExtremelyLightBlue}{RGB}{235, 245, 255}

\definecolor{myLightBlueFrame}{RGB}{97, 162, 220} 
\definecolor{mySoftBlue}{RGB}{205, 230, 255} 
\definecolor{myVeryLightBlue}{RGB}{220, 235, 255}   
\definecolor{myVeryLightYellow}{RGB}{255, 245, 204}
\definecolor{EvoPurple}{HTML}{F3F0FF}
\definecolor{myExtremelyLightBlue}{RGB}{235, 245, 255}
% partial support
\usepackage{tikz}
\usepackage{wrapfig}
\usepackage{textcomp}
\usepackage{scalerel}
\usepackage{tipa}
\usepackage{pifont}
\usepackage[subtle]{savetrees}
\usepackage[most]{tcolorbox}
\usepackage{listings}
\usepackage{caption}
\usepackage[normalem]{ulem}
\usepackage{bbm}
\usepackage{algorithmic}
\usepackage{float}

\graphicspath{{images/}}

\definecolor{HeaderGray}{gray}{0.95}
\definecolor{GroupGreen}{RGB}{220,240,220}
\definecolor{GroupBlue}{RGB}{220,230,245}
\definecolor{GroupYellow}{RGB}{250,243,210}
\definecolor{HeaderGray}{HTML}{EFEFEF}
\definecolor{GroupYellow}{HTML}{FFF2CC}
\definecolor{GroupBlue}{RGB}{91,157,218}
\definecolor{GroupOrange}{RGB}{255, 235, 200}
\newcommand{\ours}{\textsc{MMAgentMem}\xspace}
\newcommand{\sheng}[1]{{\color{blue} {\bf Sheng:} #1}}

\newcommand{\cmark}{\textcolor{green!60!black}{\ding{51}}}  % full support
\newcommand{\xmark}{\textcolor{red!75!black}{\ding{55}}}    % no support

\newcommand{\pmark}{\textcolor{GroupBlue}{\ding{51}\rotatebox[origin=c]{-6.2}{\kern-0.7em\ding{55}}}}


\newtheorem{definition}{Definition}

\title{WorldMemArena: Evaluating Multimodal Agent Memory Through Action–World Interaction}

\author{%
  Anonymous Authors \\
  Anonymous Institution \\
  \texttt{anonymous@anonymous.org} \\
}

\begin{document}
\raggedbottom

\maketitle




% =========================================================================
% PREVIOUS VERSION OF ABSTRACT (kept for reference)
% -------------------------------------------------------------------------
% \begin{abstract}
% Multimodal large language models are increasingly deployed as long-horizon agents in dynamic environments, where memory must preserve useful experience, update obsolete information, and retrieve evidence for future decisions. However, current evaluations primarily test whether systems can remember information they have seen before, rely on final answer accuracy, and often convert multimodal evidence into text, limiting their ability to assess how memory supports interaction with changing worlds over time. To address these gaps, we formulate multimodal agent memory as an \textit{Action--World Interaction Loop} and define a four stage lifecycle of \textit{writing}, \textit{maintenance}, \textit{retrieval}, and \textit{use} to diagnose where and why memory fails. Building on this formulation, we introduce \textbf{WorldMemArena}, a multimodal multi-session benchmark of 400 long-horizon tasks across two regimes: \textit{Lifelong Evolution}, tracking personal and task states that evolve across sessions, and \textit{Agentic Execution}, requiring memory over realistic observations, actions, and feedback.  Each session is annotated with gold memory points, state updates, distractors, and evidence chains to support diagnosis at each stage of the memory lifecycle. Experiments with long context agents, manually designed memory systems, and harness-based agents show that \textit{\textbf{(1)}} better memory writing and storage do not guarantee better performance; \textit{\textbf{(2)}} multimodal memory still struggles to fully use visual evidence; \textit{\textbf{(3)}} systems are unstable across domains and degrade on realistic agentic trajectories; and \textit{\textbf{(4)}} harness memory is more flexible but remains costly and less reliable.
% \end{abstract}
% =========================================================================

\begin{abstract}
Multimodal large language models are increasingly deployed as long-horizon agents, where memory must do more than recall: it must track an evolving world, revise what has gone stale, and surface the right evidence at decision time. Existing benchmarks measure recall over static dialogue, collapse memory into a single end-of-task accuracy, and reduce visual observations to captions, leaving us unable to localize failures to writing, maintenance, retrieval, or use. The rise of agent harnesses that author their own memory sharpens this gap, since we have no principled way to compare hand-designed pipelines with self-managing alternatives. To close these gaps, we formulate multimodal agent memory as an \textit{Action--World Interaction Loop} with an observable four-stage lifecycle, and instantiate it in \textbf{WorldMemArena}: 400 multi-session multimodal tasks spanning \textit{Lifelong Evolution} (evolving personal and task states) and \textit{Agentic Execution} (memory from real observations, actions, and feedback), annotated with gold memory points, updates, distractors, and evidence chains for stage-level diagnosis. This enables the first head-to-head comparison of long-context, hand-designed (RAG and external-memory), and harness-based memory agents. Results show that:
\textit{\textbf{(1)}} better memory writing and storage do not guarantee better performance; \textit{\textbf{(2)}} multimodal memory still struggles to fully use visual evidence; \textit{\textbf{(3)}} systems are unstable across domains and degrade on realistic agentic trajectories; and \textit{\textbf{(4)}} harness memory is more flexible but remains costly and less reliable.
\end{abstract}

% To address these gaps, we introduce \textbf{WorldMemArena}, a unified benchmark for multimodal agent memory with 400 multi-session multimodal interaction tasks. It formulates memory as an action and world interaction loop, and evaluates the full memory lifecycle, including writing, maintenance, retrieval, and answering. The benchmark covers \textit{Synthetic Long Horizon World} and \textit{Reconstructed Agentic World}, testing memory over evolving personal and task states as well as realistic agent trajectories. 

\begin{figure}[H]
\centering
\includegraphics[width=0.86\linewidth]{fig/test.pdf}
\vspace{-0.6em}
\caption{\textbf{Overview of WorldMemArena.} }
\label{fig:teaser}
\vspace{-0.8em}
\end{figure}


\section{Introduction}
\label{sec:intro}

% =========================================================================
% PREVIOUS VERSION OF INTRODUCTION (kept for reference)
% -------------------------------------------------------------------------
% \begin{wrapfigure}{r}{0.5\textwidth}
%   \centering
%   \includegraphics[width=\linewidth]{fig/intro.pdf}
%   \caption{An overview of multimodal agent memory evaluation.}
%   \label{fig:intro}
% \end{wrapfigure}
%
% Large language models are evolving from question answering systems into agents that operate in dynamic environments over long horizons. In this setting, memory is not merely recall over past text, but a capability for maintaining task state, tracking environmental changes, and reusing prior experience. An effective agent memory system must preserve information that remains relevant and make it available for future decisions. Thus, memory is a core capability for long horizon agents, rather than simply context extension or history retrieval. However, multimodal agent memory remains insufficiently understood and evaluated. \textbf{(i)} most existing benchmarks are still centered on static recall, and thus fail to capture sustained interaction between agents and dynamic environments, as shown in Figure~\ref{fig:intro}(a).  \textbf{(ii)} as shown in Figure~\ref{fig:intro}(b), existing evaluations often lack lifecycle level diagnosis, reducing memory ability to final question answering accuracy without assessing whether evidence is written, maintained and used correctly, which makes failure sources difficult to localize.
% \textbf{(iii)} Multimodal and agentic evidence is often converted into textual information, limiting the use of visual and environmental signals in retrieval and reasoning, see Figure~\ref{fig:intro}(c).  \textbf{(iv)} Many systems still rely on fixed human designed pipelines, and the gap between this paradigm and agent-managed memory remains underexplored.
%
% With the emergence of agent harness systems such as OpenClaw and Codex, agents are increasingly able to record experience and adapt memory use through interaction with the world. Inspired by \textit{Rich Sutton's Bitter Lesson}, we revisit the design paradigm of agent memory and ask the following question:
%
% \begin{tcolorbox}[colframe=black!50, colback=cvprblue!8, boxrule=1.5pt, arc=2mm, top=4pt, bottom=4pt, left=4pt, right=4pt,  boxsep=1pt]
% \raisebox{-0.2\baselineskip}{\includegraphics[height=1\baselineskip]{fig/idea.png}} \textit{\textbf{When agents can continuously act in realistic environments and manage their own experience, should agent memory still rely on human designed memory mechanisms?}}
% \end{tcolorbox}
%
% Motivated by this question, \textbf{(I)} we first define multimodal agent memory as an
% \textit{\textbf{Action-World Interaction Loop}}. Under this view,
% memory is not a static cache of facts, but an interaction mechanism that
% enables agents to reuse past experience in future tasks. Agents acquire
% experience through observations, actions, and feedback, while the memory
% system maintains information that remains useful for future decisions. \textbf{(II)} We introduce \textbf{WorldMemArena}, a unified benchmark for multimodal agent memory with 400 multi-session multimodal interaction tasks. It covers two settings: \textit{Synthetic Long Horizon World}, which evaluates evolving personal and task states through lifelong dialogue and long horizon projects, and \textit{Reconstructed Agentic World}, which evaluates memory over actions, feedback, and environment changes from realistic agent-environment trajectories. Under a unified framework, we assess the full memory lifecycle from writing and maintenance to retrieval and answering, as shown in Figure~\ref{fig:intro}(d).
%



% We evaluate long-context agents, human designed memory systems such as retrieval-augmented generation (RAG) and external-memory agents, as well as harness based memory agents under a unified setting. Our results show that:
% \textit{\textbf{(1)}} stronger backbones do not necessarily yield better memory performance, and memory design remains critical;
% \textit{\textbf{(2)}} memory coverage and question-answering correctness are clearly misaligned, indicating that correct storage does not imply effective use;  \textit{\textbf{(3)}} current multimodal memory still relies heavily on text captions, making visual evidence hard to use; \textit{\textbf{(4)}} systems are unstable across domains and degrade notably
% on real-environment agentic trajectories; and
% \textit{\textbf{(5)}} human designed memory systems are more structured but less adaptive, while harness based memory agents are more flexible but remain
% costly and less reliable. To sum up, our contributions are listed as follows:
%
% \begin{itemize}[left=0.2cm]
%     \item  We formulate multimodal agent memory system as an Action World Interaction Loop and define a four stage lifecycle covering writing, maintenance, retrieval, and use.
%     \item We introduce \textbf{WorldMemArena}, a unified multi-session multimodal benchmark that evaluates memory over evolving personal and task states and scenarios built from real agent trajectories.
%     \item  We evaluate long context agents, human designed memory systems, and harness based agents, showing that stored memory is not reliably used, visual evidence remains underused, systems degrade on agent scenarios, and harness based memory is flexible but costly and less reliable.
% \end{itemize}
% =========================================================================


Multimodal large language models~\cite{gpt54,qwen35_2026, claude_opus46_2026} are turning from question answering systems into agents that act in dynamic environments over long horizons~\cite{steinberger2025openclaw, claudecode2026}.  In this setting, memory is no longer simply a cache of past text, but a mechanism for tracking task state, learning from actions, and supporting decisions  through real-world interaction. A capable long horizon agent should not only recall the past, but also write useful information, revise outdated memories, and retrieve the right evidence for future decisions. How well current memory systems can fulfill this role remains insufficiently evaluated.

\begin{wrapfigure}{r}{0.45\textwidth}
  \centering
  \includegraphics[width=\linewidth]{fig/intro.pdf}
  \caption{\fontsize{9.5pt}{10pt}\selectfont Overview of multimodal agent memory evaluation.
(a) Recall-only evaluation. (b) Missing lifecycle-level diagnosis.
(c) Low-pressure chat-like settings. (d) WorldMemArena evaluates multimodal memory through action-world interaction.}
  \label{fig:intro}
  \vspace{-0.5cm}
\end{wrapfigure}

Existing benchmarks fall short of this picture in three connected ways. \textbf{\textit{(i)}} They are often built around long dialogues or extended contexts~\cite{ jiayang2026amemgyminteractivememorybenchmarking}, testing what models can remember rather than how they use past experience to guide future actions.(Figure~\ref{fig:intro}(a)).  \textbf{\textit{(ii)}} As shown in Figure~\ref{fig:intro}(b), many evaluations~\cite{zhao2026amabenchevaluatinglonghorizonmemory, hu2026evaluatingmemoryllmagents} report only final question answering accuracy, without checking whether relevant evidence is written, updated, retrieved, and used at the right time, making it difficult to identify where memory failures occur. \textbf{\textit{(iii)}} Figure~\ref{fig:intro}(c) shows that existing benchmarks remain largely text-centric, often converting images into captions before evaluation, with limited real interaction and insufficient pressure on multimodal evidence use.

Beyond these evaluation limitations, current benchmarks also miss a deeper shift in how agent memory is built and used. Agent harness systems such as OpenClaw~\cite{steinberger2025openclaw} and Codex~\cite{codex2026} now let agents author and reorganize their own memory during interaction, blurring the line between the memory module and the policy that uses it. In the spirit of \textit{Sutton's Bitter Lesson}, this invites a question the field should be asking head-on:
% Beneath these methodological gaps lies a more fundamental shift that current evaluations have yet to confront. Agent harness systems such as OpenClaw and Codex now let agents author and reorganize their own memory during interaction, blurring the line between the memory module and the policy that uses it. In the spirit of \textit{Sutton's Bitter Lesson}, this invites a question the field should be asking head-on:

\vspace{-0.12cm}
\begin{tcolorbox}[colframe=black!50, colback=cvprblue!8, boxrule=1.5pt, arc=2mm, top=2pt, bottom=2pt, left=4pt, right=4pt,  boxsep=1pt]
\raisebox{-0.2\baselineskip}{\includegraphics[height=1\baselineskip]{fig/idea.png}} \textit{\fontsize{9.3pt}{10.2pt}\selectfont When agents can continuously act in realistic environments and manage their own experience, should agent memory be evaluated as a predefined write-update-retrieve pipeline, or as a capability formed through interaction and used to support future decisions over time?}
\end{tcolorbox}

Answering this question requires an evaluation that treats memory as a process rather than a static snapshot. We reframe multimodal agent memory as an \textit{\textbf{Action and World Interaction Loop}}. At each step, the agent observes a partially visible world, takes an action, receives feedback, and uses memory to guide future actions and retain useful evidence. Under this view, memory has an observable lifecycle that covers what is written, how it is maintained as the world changes, what evidence is retrieved, and how the retrieved evidence is used. As shown in Figure~\ref{fig:intro}(c), each stage can be evaluated using shared trajectory evidence, rather than inferred from a single accuracy score.

% Answering this question requires an evaluation that treats memory as a process, not a snapshot. We therefore reframe multimodal agent memory as an \textit{\textbf{Action–World Interaction Loop}}: at each step the agent observes a partially visible world, takes an action, receives feedback, and uses memory both to choose its action and to decide what about this step is worth keeping. Under this view memory has an observable lifecycle — what is written, how it is maintained as the world changes, what is retrieved, and how retrieved evidence is finally used — and each phase can be probed with shared trajectory evidence rather than inferred from a single accuracy number (Figure~\ref{fig:intro}(d)).

We instantiate this view in \textbf{WorldMemArena}, a multimodal multi-session benchmark of 400 long-horizon interaction tasks spanning two complementary regimes. \textit{Lifelong Evolution} focuses on personal and task states that evolve across sessions, requiring systems to continuously track, update, and reuse long-term memories. \textit{Agentic Execution} places memory in realistic agent trajectories, where systems must extract reusable evidence from observations, actions, and feedback rather than relying on pre-organized textual narratives. Each session is annotated with gold memory points, state updates, distractors, and answer supporting evidence chains. These annotations support diagnosis across memory writing, maintenance, retrieval, and use, while providing a shared evidence base for comparing different memory systems.

Under a unified setting, the evaluation covers long-context agents, manually designed memory systems, and memory agents built on execution harnesses. The results reveal four findings: 
\textit{\textbf{(1)}} high memory coverage \sheng{what is memory coverage} does not guarantee correct memory use;
\textit{\textbf{(2)}} multimodal  remains a major bottleneck, especially for complex visual reasoning tasks;
\textit{\textbf{(3)}} memory performance varies across domains and drops on realistic agentic trajectories\sheng{what is this realistic agentic trajectories, how it relates to different domains?}; and
\textit{\textbf{(4)}} manually designed memory systems are more structured but less adaptive, while harness based memory agents are more flexible but remain costly and less reliable. To sum up, our contributions are listed as follows:

\begin{itemize}[left=0.2cm]
    \item We formulate multimodal agent memory as an \textit{Action--World Interaction Loop} and define a four stage lifecycle of writing, maintenance, retrieval, and use.

    \item We introduce \textbf{WorldMemArena}, a multi-session multimodal benchmark covering \textit{Lifelong Evolution} and \textit{Agentic Execution}, with annotations for stage level memory diagnosis.

    \item We conduct a unified comparison of three representative agent memory paradigms, identifying their respective strengths, failure modes, and implications for future design.
\end{itemize}

% We instantiate this view in \textbf{WorldMemArena}, a multimodal multi-session benchmark of 400 long-horizon interaction tasks spanning two complementary regimes. The \textit{Synthetic Long-Horizon World} simulates evolving personal lives and multi-stage projects whose hidden state must be tracked across sessions, isolating the demands of persistent state evolution. The \textit{Reconstructed Agentic World} replays real GUI and embodied trajectories so that memory must be formed from actions, screenshots, and environment feedback rather than from narration. Every session is annotated with gold memory points, updates, and interference items, letting us score writing, maintenance, retrieval, and use under a single protocol — and, critically, place long-context models, hand-designed memory systems (RAG and external-memory agents), and harness-based memory agents on the same axis for the first time.

% The picture that emerges is consistent and somewhat uncomfortable. Stronger backbones and larger memory stores do not translate into better answers: what is stored is routinely not what is used, retrieval rather than capacity is the dominant bottleneck, and visual evidence remains hard to reuse because systems still route it through text. Performance also degrades sharply when we move from synthetic dialogues to real agent trajectories, exposing how brittle today's designs are to the conditions they will actually face. Harness-based agents that manage their own memory do gain flexibility on these harder regimes, but at a cost in stability and overhead — suggesting the path forward is neither pure hand-crafted pipelines nor fully unconstrained self-management, but designs that combine the structure of the former with the adaptivity of the latter.
% \textbf{Contributions.} We make three contributions: \textit{(i)} we formulate multimodal agent memory as an Action–World Interaction Loop with a four-stage observable lifecycle covering writing, maintenance, retrieval, and use; \textit{(ii)} we release \textbf{WorldMemArena}, a multi-session multimodal benchmark spanning evolving long-horizon worlds and reconstructed agent trajectories, with annotations that enable lifecycle-level diagnosis; and \textit{(iii)} we provide the first head-to-head comparison of long-context, human-designed, and harness-based memory under a unified protocol, showing where each paradigm succeeds, where it fails, and what this implies for the next generation of agent memory systems.



\section{Related Works}
\label{sec:related}

\textbf{Memory Benchmarks and Evaluation.} Early memory benchmarks such as LoCoMo~\cite{maharana2024evaluating}, MemoryAgentBench~\cite{hu2025evaluating}, and Realme~\cite{bian2026realmembenchmarkingllmsrealworld} focus on long-dialogue settings, measuring whether models can retain and recall historical information. These benchmarks treat memory as static recall over text and do not capture how memory supports dynamic task execution. More recent agent-oriented benchmarks~\cite{he2026memoryarenabenchmarkingagentmemory, zhao2026amabenchevaluatinglonghorizonmemory} incorporate tool traces, environment feedback, and task dependencies, moving closer to realistic agent-environment interaction. However, evaluation still centers on final success rates or question answering accuracy, making it difficult to identify where and why memory fails. WorldMemArena differs by decomposing evaluation into writing, maintenance, retrieval, and use, making it possible to localize where memory failures originate.

\textbf{Multimodal Memory Mechanisms.} Recent multimodal memory systems~\cite{long2025seeinglisteningrememberingreasoning, liu2025memversemultimodalmemorylifelong, zhou2026videomemoryconsistentvideogeneration, fu2026latentmemcustomizinglatentmemory} have demonstrated strong capabilities in visual understanding and long-term information retention. Their evaluations, however, are largely confined to image and video comprehension tasks, with limited attention to how memory operates within agent interaction loops. Benchmarks that incorporate multimodal memory~\cite{bei2026memgallerybenchmarkingmultimodallongterm, lu2026mmamultimodalmemoryagent} extend evaluation to images, videos, and dialogues, but cover a narrow range of scenarios and apply limited evaluation pressure on evidence reuse. WorldMemArena broadens the scope to multi-session agent interaction, testing whether systems can preserve, update, and reuse multimodal evidence as tasks and environments evolve.


\section{Problem Formulation} 


\subsection{Memory as an Action-World Interaction Loop}
\label{sec:problem}
We define each instance as a long horizon agent-world interaction process. Given an initial task context $x$, the agent does not directly observe the full world state. At step $t$, the world has a latent state $z_t$, from which the agent receives an observation $o_t$. The agent then selects an action $a_t$ based on the observation and its current memory state $m_t$. After the action is executed, the environment updates its state and returns feedback $f_t$:
\vspace{-0.3em}
\[
o_t=\Omega(z_t), \qquad a_t=\pi(o_t,m_t), \qquad (z_{t+1},f_t)=\mathcal{E}(z_t,a_t).
\]
Here, $\Omega$ maps the latent world state to observable inputs, $\pi$ denotes the agent policy, and $\mathcal{E}$ represents the environment response, including both state transition and feedback generation. Observations may include  language, visual inputs or logs, while actions may include responses, tool calls, or execution.

Based on the above process, we denote the full trajectory as
$\tau=(x;\eta_1,\ldots,\eta_T)$, where each event
$\eta_t=(o_t,a_t,f_t)$ records the observation, action, and feedback at step
$t$. To evaluate long-horizon memory, we further segment the trajectory into
sessions, i.e.,
$\tau=\tau^{(1)}\circ\tau^{(2)}\circ\cdots\circ\tau^{(S)}$.
Within each session, the agent only observes local context, while the world state persists and evolves across sessions. This creates a natural point: later decisions may depend on evidence that is no longer directly visible, and we focus on whether the agent can recover and use such evidence through memory.


\subsection{Memory Lifecycle as a Diagnostic Framework}
\label{sec:problem-lifecycle}
The Action World Interaction Loop in \S\ref{sec:problem-lifecycle} is architecture agnostic. It does not assume where memory is stored or how it is represented. This allows us to evaluate different memory systems through four observable phases of writing, maintenance, retrieval, and use. These phases capture the shared lifecycle of preserving and reusing information across sessions.

\textbf{$\spadesuit$ Observe to Write.} This phase evaluates  \textit{whether the system can identify future useful evidence
from the current session.} Given the previous memory state $m_{s-1}$ and the
current session trajectory $\tau^{(s)}$, the system produces a memory delta
$\Delta_s=\textsc{Write}(m_{s-1},\tau^{(s)})$. The objective is selective
retention, keeping information that may support future responses or actions
rather than storing the full trajectory.

\textbf{$\spadesuit$ Update and Consolidate.} This phase evaluates \textit{how newly written information is integrated into
existing memory.} The system updates its state as
$m_s=\textsc{Maintain}(m_{s-1},\Delta_s)$.  Since long-horizon interaction is
not purely additive, memory must support revision and consolidation as user
preferences, task states, and environmental evidence evolve.

\textbf{$\spadesuit$ Retrieve for Decision.} This phase evaluates \textit{whether the system can access the right evidence when a
future query or decision need arises.} For a query $q$, retrieval returns
$R_{s,q}=\textsc{Retrieve}(m_s,q)$. The goal extends beyond semantic similarity to decision relevance, requiring the retrieved context to contain evidence needed for the current answer or action.

\textbf{$\spadesuit$ Use and Act.} This phase evaluates \textit{whether retrieved memory is faithfully used in the final
response or action.} Given $q$ and retrieved evidence $R_{s,q}$, the system
outputs $\hat{y}_{s,q}=\textsc{Answer}(q,R_{s,q})$. Failures may still arise
when the system ignores relevant evidence, relies on outdated memory, or fails
to translate prior experience into appropriate action.


\input{table/data}

\section{\textit{WorldMemArena}: Agent Memory in Action-Wolrd Interaction}
\label{sec:benchmark}

\textbf{Overview.} WorldMemArena consists of 400 multi-session multimodal interaction tasks across two regimes (\textit{Lifelong Evolution} and \textit{Agentic Execution}). Each task is a temporally ordered sequence of sessions, where the agent receives partial observations and must rely on memory to inform decisions in later sessions. To support fine-grained diagnosis, every session is annotated with three types of structured labels. \textit{Gold memory points} specify the information that should be retained after a session, representing ground-truth memory content. \textit{State updates} mark where previously stored information becomes outdated and must be revised, testing whether the memory system can maintain temporal consistency. \textit{Distractors} introduce plausible but irrelevant or superseded information, testing whether the system can distinguish currently valid evidence from noise. In addition, each question is paired with \textit{evidence points}, the subset of gold memory points that are necessary to answer it correctly. These annotations together enable evaluation at each stage of the memory lifecycle.


\subsection{Memory Regimes}
\label{sec:regimes}

% \textbf{Agentic Execution.} This task models memory demands in realistic agent execution. Each instance is derived from a real or realistic agent trajectory, capturing what the agent observes, how it acts, and how the environment responds. Within each task, later steps depend on earlier outcomes and experience. The goal is to evaluate whether systems can form reusable agentic memory about past actions, state changes, and execution lessons that support future decisions and actions.
\textbf{Agentic Execution.} Each instance is derived from a real or realistic agent trajectory containing observations, actions, and environment feedback. Later steps depend on earlier outcomes, so the agent must convert past execution experience into reusable memory that informs future decisions.

\textbf{Lifelong Evolution.} Each instance is generated from a hidden world state that evolves across sessions. It covers two scenarios: (1) \textit{lifelong personal evolution}, where scattered interactions must be consolidated into coherent personal memory; and (2) \textit{long-horizon projects}, where task goals, intermediate results, and feedback shift across stages, requiring the agent to maintain up-to-date progress memory.

% \textbf{Lifelong Evolution.}  This task simulates evolving personal and task states in realistic long-term settings. Each instance is generated from a hidden world state, allowing earlier experiences, decisions, and intermediate progress to affect later interactions. It covers two scenarios: (1) \textit{lifelong evolution}, where the system must consolidate scattered interactions into reusable personal memory for future responses and actions; and (2) \textit{long-horizon projects}, which simulate how task goals, intermediate results, and feedback change across multiple project stages, requiring the system to maintain an up-to-date memory of progress to support future decisions and actions.

\textbf{Why both Regimes are Needed.} As the Action-World Interaction Loop requires the agent to both observe an evolving world and act within it, two demands on memory arise:
\textit{(1) Persistent state tracking} requires maintaining an accurate representation of an evolving world across sessions, which is evaluated by \textit{Lifelong Evolution} through controlled state evolution.
\textit{(2) Action grounded experience reuse} requires turning observations, action outcomes, and feedback into knowledge for later decisions, which is evaluated by \textit{Agentic Execution} through realistic execution trajectories.


\begin{figure}[t]
\centering
\includegraphics[width=\linewidth]{fig/pipline.pdf}
\caption{\fontsize{9.5pt}{10pt}\selectfont Data construction pipeline and benchmark composition. (a) WorldMemArena constructs data around two task regimes, \textit{Lifelong Evolution} and \textit{Agentic Execution}, by segmenting sessions, extracting and updating gold memory points, removing redundancy, and constructing midway and final QA checkpoints. (b-1) The benchmark covers both GUI and embodied interaction settings. (b-2) The upper charts summarize gold memory points across the benchmark, including update memory points and distractor memory points. The lower chart shows the task distribution across domains in \textit{Lifelong Evolution}.}
\label{fig:pipeline}
\vspace{-0.4cm}
\end{figure}

\vspace{-0.1cm}
\subsection{Data Collection}
\vspace{-0.2cm}
As shown in Figure~\ref{fig:pipeline}(a), WorldMemArena is constructed through a unified automated memory construction pipeline with four steps.
\textit{{\textbf{(1)}}} Raw data is segmented into multi-session instances. For \textit{Lifelong Evolution}, a hidden world state is first defined and sessions are generated in temporal order, each revealing partial information about a persona or project. For \textit{Agentic Execution}, existing agent trajectories are split at subgoal boundaries, key feedback points, or state changes. \textit{\textbf{(2)}} For each session window, gold memory points are extracted, covering facts to retain, state updates to revise, and evidence required by future questions. \textit{\textbf{(3)}} Memory points are merged, revised, and deduplicated across sessions to remove redundancy and ensure temporal consistency.  \textit{\textbf{(4)}} Question-answer pairs are constructed from the refined gold memory points, covering 11 question types. Each instance is further reviewed by 2-3 human annotators to ensure quality.


\subsection{\textcolor{black}{Data Statics}}

\textbf{Dataset Scale and Coverage.} Table~\ref{tab:bench-cmp-new} compares WorldMemArena with existing benchmarks. Prior datasets typically focus on either long-form dialogue or agentic trajectories, whereas this benchmark covers both lifelong evolution and agentic execution. It contains 400 multi-session samples, with an average of 18.4 sessions and approximately 9.1K tokens per sample, making it substantially longer than existing multimodal memory benchmarks. It further provides 24{,}258 QA pairs and 15{,}595 images or screenshots, supporting broader question coverage and richer visual grounding. Most existing benchmarks do not evaluate the full memory lifecycle; the closest prior work, HaluMem, addresses memory storage and recall but remains limited to the textual modality.

\textbf{Domain and Annotations.} As shown in Figure~\ref{fig:pipeline}(b), \textit{Lifelong Evolution} covers 6 domain specific project types, with each session containing an average of 4 images and 15-20 dialogue turns. \textit{Agentic Execution} preserves real agent execution traces and their corresponding visual states, covering 6 GUI subcategories and 4 Embodied subcategories. Across both regimes, fine-grained lifecycle annotations are provided.
Each session contains an average of 10 key memory points, 3 update points,
and 2 interference points. Each sample further includes staged QA
checkpoints with an average of 5 evaluation positions. Each question is
paired with retrieval evidence, where most require 1-2 evidence items
and more complex questions require 5-6, covering both textual and visual information.


\subsection{Evaluation Protocol}
Following the four lifecycle stages defined in \S\ref{sec:problem-lifecycle}, we evaluate whether a memory system can correctly write, maintain, retrieve, and use memory across long-horizon interactions. 

\textit{\textbf{Stage 1.}} For each session, newly written memory items are matched against the gold memory points introduced in that session, with \textit{memory recall} used as the coverage metric. Each written item is further assessed by an LLM-as-a-Judge and classified as \textit{correct, hallucinated, or irrelevant}, distinguishing effective memory writing from noisy or unsupported storage.

\textit{\textbf{Stage 2.}} For gold memory points marked as updates, the system memory after the corresponding session is examined to determine whether the new information is preserved and the obsolete version is properly handled. \textit{An update is considered successful only when the revised memory is retained and the old version is removed or overwritten}. This criterion prevents simple accumulation of historical information from being misclassified as effective memory maintenance.

\textit{\textbf{Stage 3.}} For each checkpoint question, the retrieved memory items are matched against the annotated gold evidence. The evidence may be grounded in either textual or visual information, and all evidence types are evaluated under a unified coverage criterion. \textit{Recall} measures whether the required evidence is retrieved, while \textit{NDCG} measures whether relevant evidence is ranked near the top, thereby separating retrieval quality from final answer correctness.

\textit{\textbf{Stage 4.}} Checkpoint questions are grouped into four categories and twelve capability axes: \textit{Basic} covers factual recall; \textit{Robustness} covers dynamic update, memory boundary, and memory conflict; \textit{Reasoning} covers temporal reasoning, knowledge reasoning, and test-time learning; and \textit{Multimodal} covers visual fact recall, visual search, visual update, and cross-modal reasoning. Each question is jointly evaluated using LLM-as-a-Judge, F1, and BLEU to reduce biases from any single metric.

% Implementation details are provided in \textcolor{red}{Appendix C}.

% \input{table/table2}


\section{Experiments}
\label{sec:exp}

\subsection{Setup}
We evaluate three mainstream memory paradigms on \textbf{WorldMemArena}. 

\textbf{Long-Context Agents.} To test whether frontier models can handle long-horizon memory tasks by relying solely on context, these agents concatenate the full interaction history into the prompt as in-context memory, without explicit abstraction, updating, or retrieval. We evaluate GPT-5.4-mini, Qwen3.5 plus, Gemini 3 flash, DeepSeek V4 and Claude Haiku 4.5. As no independent memory state is exposed, only final question-answering performance is measured.

\textbf{Manually  Designed Memory Systems.} To assess whether explicitly engineered memory mechanisms can improve memory construction, maintenance, retrieval, and downstream use, we evaluate two types of systems. \textit{External memory agents} such as MemGPT and Mem0 perform information abstraction, consolidation, and retrieval through learned or hand-crafted modules. \textit{RAG systems} such as UniversalRAG  store historical information in an indexed document store and access it via retrieval. To control for backbone differences, all systems use GPT-5.4-nano as the base model. Because these systems expose observable memory states and retrieval outputs, the full memory lifecycle can be evaluated.

\textbf{Harness-Based Memory Agents.} To examine whether agents can autonomously manage memory without a fixed external module, we evaluate agent harnesses where memory is written, maintained, retrieved, and used by the harness itself during interaction. We test OpenClaw paired with GPT-5.4 and DeepSeek-V4, and Codex paired with GPT-5.4, feeding session contexts sequentially and testing with staged checkpoint QA. Since the internal memory process is difficult to decompose, we primarily conduct end-to-end evaluation.

\input{table/table1}

\subsection{\textcolor{black}{Main Results}}
Table~\ref{tab:data150-gpt-quality} reports the overall performance of different human designed systems across the full memory lifecycle. We identify four main findings.\ding{182} \textbf{Multimodal memory is still not effectively used.}  Text-based systems such as MemoryGPT and A-Mem achieve more stable final answer quality, while multimodal systems such as ViLoMem and MIRIX show limited downstream gains despite access to visual inputs. This suggests that current systems still struggle to encode and reuse visual evidence as reliable long term memory. \ding{183} \textbf{High memory quality does not necessarily lead to high QA quality.} High memory quality does not necessarily lead to high QA quality. Qwen3-VL-Embedding and M2A perform well in memory storage and recall, but their final answers remain limited. This indicates that correct memory writing is insufficient; systems must also retrieve and use the right evidence during answer generation. \ding{184} \textbf{Retrieval remains a key bottleneck for final performance.} MemoryGPT achieves the strongest evidence retrieval and answer correctness, while A-Mem uses retrieved information effectively despite lower memory coverage. In contrast, AUGUSTUS constructs reasonably good memories but fails to surface key evidence at inference time, limiting its final QA performance. \ding{185} \textbf{Most systems remain weak in memory updating and distractor rejection.}  Nearly all systems are brittle under information changes and interfering content, indicating that they tend to accumulate memories rather than maintain a consistent long-term state. This suggests that current human designed memory systems still focus more on how much they remember than on how well they maintain and update memory over time.

\input{table/table3}
Table~\ref{base} compares final answer performance between long context agents and harness based memory agents. Most long context agents perform poorly, with some falling below dedicated memory systems, indicating that the benchmark requires long horizon evidence integration rather than context extension alone. DeepSeek V4 benefits mainly from its larger context window, while standard context models remain limited. Harness based memory agents outperform most human designed memory systems, suggesting that agent managed memory is more flexible. However, the same backbone performs differently across harnesses, showing that native memory design and adaptation mechanisms also affect final performance.


\section{\textcolor{black}{Analysis}}

\noindent \textbf{[RQ1] \textit{Where do memory failures occur in the lifecycle?}}

\begin{wrapfigure}{r}{0.43\textwidth} 
\vspace{-5em}
  \centering
  \includegraphics[width=\linewidth]{fig_RQ/RQ1.pdf}
  \vspace{-1.5em}
\caption{\fontsize{9.5pt}{10pt}\selectfont 
(a) shows the trend of average QA accuracy across dialogue sessions.
(b), (c) summarize memory-point composition in terms of update operations (Append, Revise, Delete, and Merge) and fact salience.
}
  \label{trade1}
  \vspace{-1.5em}
\end{wrapfigure}

\textit{Memory failures occur across the full lifecycle and compound over time.} \textit{(i)} Figure~\ref{trade1}(a) shows that storing more memories does not necessarily make them usable; even with high storage coverage, systems may fail to retrieve the key evidence needed for the current decision. \textit{(ii)} As illustrated in Figure~\ref{trade1} (b), most systems rely on append only updates, adding new information when evidence changes rather than revising, removing, or reorganizing obsolete memories. \textit{(iii)} Over long trajectories, Figure~\ref{trade1}(c) captures a compounding pattern in which early omissions reduce later evidence availability, while incorrect outputs may contaminate future memory updates and further induce hallucinated answers.


\noindent  \textbf{[RQ2] \textit{Are memory system designs constrained by domain-specific data?}}

\textit{Memory performance varies across domains.} As shown in Figure~\ref{fig:placeholder}(a-b), most systems perform better in \textit{Lifelong Evolution} than in \textit{Agentic Execution}. This suggests that existing methods are more suited to explicit long-term state evolution, while extracting usable memory from action traces and environment feedback remains challenging. Performance also differs across tasks, with long-horizon embodied tasks such as visual navigation posing greater challenges, suggesting that current systems still struggle to track memory across sessions and use it for later decisions.
 
\noindent \textbf{[RQ3] \textit{How does multimodality affect the memory lifecycle?}}

\textit{Memory systems still struggle with complex visual memory tasks.} As shown in Figure~\ref{fig:placeholder}(c), systems perform relatively stably on simple visual fact recall, but degrade on tasks that depend on long interaction histories, such as cross-modal reasoning. This suggests that the core challenge of multimodal memory is to maintain visual states over time and integrate visual evidence with historical context when needed.

\begin{figure}[h]
\vspace{-0.3em}
    \centering
    \includegraphics[width=1\linewidth]{fig_RQ/RQ2.pdf}
    \vspace{-1.5em}
\caption{\small
Performance comparison across scenarios and visual QA tasks.
(a) Heatmap of baseline performance across fine-grained task categories in the agent-scenario and long-dialogue settings.
(b) Average baseline performance under the two settings, where long-dialogue tasks show higher overall performance than agent-scenario tasks.(c) Box plots of baseline performance across three visual QA task types.
}
    \label{fig:placeholder}
    % \vspace{-1em}
\end{figure}


\noindent \textbf{[RQ4]  \textit{What strengths and limitations do different memory systems exhibit across task types?}}

\begin{wrapfigure}{r}{0.5\textwidth} 
  \centering
  \includegraphics[width=\linewidth]{fig_RQ/RQ4.pdf}
  \vspace{-1.5em}
\caption{\small
(a) Fine-grained QA performance of different baselines on individual tasks.
(b) Recall@K and NDCG@K (Normalized Discounted Cumulative Gain) trends under different retrieval cutoffs.
}
  \label{backbone}
\end{wrapfigure}

\textit{Memory performance depends more on system design than on backbone scale or retrieval volume} As shown in Table~\ref{tab:data150-gpt-quality}, most systems achieve high memory storage recall and writing quality, yet their evidence recall at question-answering time drops substantially, indicating that correctly stored memories are not effectively surfaced when needed. Figure~\ref{backbone}(b) further shows that increasing the retrieval scope does not always improve answer quality, as longer contexts may introduce redundant, outdated, or irrelevant evidence. This issue is more evident in multimodal tasks, where long interactions create substantial visual redundancy and make key visual evidence harder to locate and use.


\noindent \textbf{[RQ5] \textit{Can agents turn memory into action?}}

\textit{Past memory is not reliably converted into reusable knowledge, and experiential evidence remains fragile.} As shown in \textcolor{black}{Figure~\ref{backbone}(a)}, systems perform worse on reasoning and test-time learning tasks, suggesting that they are better at storing past information than using it to guide future decisions. Analysis of retrieved memory points in Figure~\ref{trade1}(c) shows that retrieved memories are dominated by explicit textual facts, whereas tool feedback, failed actions, and visual details are often omitted.


\noindent \textbf{[RQ6]  \textit{How far do human designed memory systems fall short in agentic memory?}}

\begin{wrapfigure}{r}{0.4\textwidth} 
  \centering
  \vspace{-1.5em}
  \includegraphics[width=\linewidth]{fig_RQ/RQ5.pdf}
  \vspace{-1.5em}
\caption{\small
Token efficiency \& QA performance trade-off among different baselines.
Circle size indicates average inference time, with larger circles denoting higher time cost.
}
  \label{harness}
    \vspace{-3em}
\end{wrapfigure}

\textit{Fixed memory architectures struggle to adapt to dynamic memory demands.} As shown in Figure~\ref{harness}(a), human designed memory systems perform comparably to harness based methods on simpler long-horizon tasks, but their fixed pipelines become limiting in complex agentic settings where memory must adapt to task feedback and environmental changes. harness based agentic memory managers are more flexible because they can record, retrieve, and revise memory during interaction. However, this result also shows that current harness-based memory remains computationally expensive and framework-dependent, limiting its stability and transferability.

\section{Discussion: Paths Forward for Multimodal Agent Memory}
\label{sec:discussion}

\textbf{\ding{118}  Memory should be treated as a controllable state, not a retrieved context.} 
Long-term memory should move beyond storage and retrieval toward continual state maintenance. Rather than treating memory as passive historical context, future systems should model it as an evolving representation of users, tasks, and environments. This requires tracking validity, version changes, and conflicts, so outdated memories can be revised, replaced, or invalidated.

\textbf{\ding{118}  Multimodal memory should transform observations into reusable interaction evidence.} 
The key challenge for multimodal memory is not storing more visual inputs, but extracting evidence that matters for future decisions. In interactive agent settings, visual observations, interface states, tool feedback, and environmental changes are often dispersed across long trajectories. Future systems should preserve links between raw observations, structured states, and semantic memories, enabling key multimodal evidence to be located and reused.

\textbf{\ding{118}  Agent memory should transform past experience into future action.} Agent memory should go beyond supporting past oriented question answering and instead contribute to future behavioral improvement. An effective long-term memory system should convert past interactions into reusable decision knowledge, including failed actions, tool feedback, environmental changes, and signals for strategy adjustment. Future benchmarks and system designs should move toward memory to action evaluation, assessing whether memory helps agents adapt strategies, avoid repeated mistakes, and transfer experience to new tasks.

\section{Conclusion}
\label{sec:conclusion}


% % ------------------------------
% \newpage

% \begin{figure}
%     \centering
% \includegraphics[width=0.5\linewidth]{fig/memory_point.pdf}
%     \caption{Fig 4}
%     \label{fig:placeholder}
% \end{figure}

% \begin{figure}
%     \centering
%     \includegraphics[width=1\linewidth]{fig/matrix.pdf}
%     \caption{Fig 5-6}
%     \label{fig:placeholder}
% \end{figure}

% \begin{figure}
%     \centering
%     \includegraphics[width=1\linewidth]{fig/RQ4.pdf}
%     \caption{Fig 7-8}
%     \label{fig:placeholder}
% \end{figure}






\clearpage
\bibliographystyle{plain}
\bibliography{neurips_2025}

\medskip

%%%%%%%%%%%%%%%%%%%%%%%%%%%%%%%%%%%%%%%%%%%%%%%%%%%%%%%%%%%%

\appendix

\section{Additional dataset details}
\label{app:data}

\subsection{Persona seed bank}
Domain A persona seeds span 24 occupational and 16 life-stage archetypes;
Domain B persona seeds span 18 project archetypes (research, product, ops,
creative). Each seed is expanded into a structured persona document
(name, role, goals, recurring stakeholders, work style) that conditions
all downstream stages.

\subsection{Quality validation}
Each generated session passes through automatic validators (memory point
coverage, image caption coverage, interference detectability, update
chain consistency) before being assembled into the dataset. Samples
failing any validator are regenerated up to 3 times.

\section{Adapter interface}
\label{app:adapter}
Every memory system implements the seven-method \texttt{MemoryAdapter}
interface: \texttt{reset}, \texttt{ingest\_turn}, \texttt{end\_session},
\texttt{snapshot\_memories}, \texttt{export\_memory\_delta},
\texttt{retrieve}, \texttt{get\_capabilities}. This unifies systems
written in Python, hosted via local servers (Qdrant, Neo4j), or wrapped
from external repositories.

\section{Judge prompts}
\label{app:judges}
We use 6 LLM-judge prompts: \emph{recall} (batch), \emph{correctness}
(batch, 3-way classification), \emph{update handling} (per-update),
\emph{interference rejection} (per-interference), \emph{answer quality}
(3-way), and \emph{retrieval coverage}. Each prompt is run with three
seeds and majority-voted.

\section{Reproducibility}
\label{app:repro}
The benchmark, eval framework, generation pipeline, and adapter wrappers
are released under an open-source license. End-to-end reproduction of
the main tables takes approximately 2{,}000 GPU-hours on Qwen-32B and
20{,}000 judge-LLM calls per memory system.

\newpage
%%%%%%%%%%%%%%%%%%%%%%%%%%%%%%%%%%%%%%%%%%%%%%%%%%%%%%%%%%%%

\section*{NeurIPS Paper Checklist}

%%% BEGIN INSTRUCTIONS %%%
The checklist is designed to encourage best practices for responsible machine learning research, addressing issues of reproducibility, transparency, research ethics, and societal impact. Do not remove the checklist: {\bf The papers not including the checklist will be desk rejected.} The checklist should follow the references and follow the (optional) supplemental material.  The checklist does NOT count towards the page
limit. 

Please read the checklist guidelines carefully for information on how to answer these questions. For each question in the checklist:
\begin{itemize}
    \item You should answer \answerYes{}, \answerNo{}, or \answerNA{}.
    \item \answerNA{} means either that the question is Not Applicable for that particular paper or the relevant information is Not Available.
    \item Please provide a short (1--2 sentence) justification right after your answer (even for \answerNA). 
   % \item {\bf The papers not including the checklist will be desk rejected.}
\end{itemize}

{\bf The checklist answers are an integral part of your paper submission.} They are visible to the reviewers, area chairs, senior area chairs, and ethics reviewers. You will also be asked to include it (after eventual revisions) with the final version of your paper, and its final version will be published with the paper.

The reviewers of your paper will be asked to use the checklist as one of the factors in their evaluation. While \answerYes{} is generally preferable to \answerNo{}, it is perfectly acceptable to answer \answerNo{} provided a proper justification is given (e.g., error bars are not reported because it would be too computationally expensive'' or ``we were unable to find the license for the dataset we used''). In general, answering \answerNo{} or \answerNA{} is not grounds for rejection. While the questions are phrased in a binary way, we acknowledge that the true answer is often more nuanced, so please just use your best judgment and write a justification to elaborate. All supporting evidence can appear either in the main paper or the supplemental material, provided in appendix. If you answer \answerYes{} to a question, in the justification please point to the section(s) where related material for the question can be found.

IMPORTANT, please:
\begin{itemize}
    \item {\bf Delete this instruction block, but keep the section heading ``NeurIPS Paper Checklist"},
    \item  {\bf Keep the checklist subsection headings, questions/answers and guidelines below.}
    \item {\bf Do not modify the questions and only use the provided macros for your answers}.
\end{itemize} 
 

%%% END INSTRUCTIONS %%%


\begin{enumerate}

\item {\bf Claims}
    \item[] Question: Do the main claims made in the abstract and introduction accurately reflect the paper's contributions and scope?
    \item[] Answer: \answerTODO{} % Replace by \answerYes{}, \answerNo{}, or \answerNA{}.
    \item[] Justification: \justificationTODO{}
    \item[] Guidelines:
    \begin{itemize}
        \item The answer \answerNA{} means that the abstract and introduction do not include the claims made in the paper.
        \item The abstract and/or introduction should clearly state the claims made, including the contributions made in the paper and important assumptions and limitations. A \answerNo{} or \answerNA{} answer to this question will not be perceived well by the reviewers. 
        \item The claims made should match theoretical and experimental results, and reflect how much the results can be expected to generalize to other settings. 
        \item It is fine to include aspirational goals as motivation as long as it is clear that these goals are not attained by the paper. 
    \end{itemize}

\item {\bf Limitations}
    \item[] Question: Does the paper discuss the limitations of the work performed by the authors?
    \item[] Answer: \answerTODO{} % Replace by \answerYes{}, \answerNo{}, or \answerNA{}.
    \item[] Justification: \justificationTODO{}
    \item[] Guidelines:
    \begin{itemize}
        \item The answer \answerNA{} means that the paper has no limitation while the answer \answerNo{} means that the paper has limitations, but those are not discussed in the paper. 
        \item The authors are encouraged to create a separate ``Limitations'' section in their paper.
        \item The paper should point out any strong assumptions and how robust the results are to violations of these assumptions (e.g., independence assumptions, noiseless settings, model well-specification, asymptotic approximations only holding locally). The authors should reflect on how these assumptions might be violated in practice and what the implications would be.
        \item The authors should reflect on the scope of the claims made, e.g., if the approach was only tested on a few datasets or with a few runs. In general, empirical results often depend on implicit assumptions, which should be articulated.
        \item The authors should reflect on the factors that influence the performance of the approach. For example, a facial recognition algorithm may perform poorly when image resolution is low or images are taken in low lighting. Or a speech-to-text system might not be used reliably to provide closed captions for online lectures because it fails to handle technical jargon.
        \item The authors should discuss the computational efficiency of the proposed algorithms and how they scale with dataset size.
        \item If applicable, the authors should discuss possible limitations of their approach to address problems of privacy and fairness.
        \item While the authors might fear that complete honesty about limitations might be used by reviewers as grounds for rejection, a worse outcome might be that reviewers discover limitations that aren't acknowledged in the paper. The authors should use their best judgment and recognize that individual actions in favor of transparency play an important role in developing norms that preserve the integrity of the community. Reviewers will be specifically instructed to not penalize honesty concerning limitations.
    \end{itemize}

\item {\bf Theory assumptions and proofs}
    \item[] Question: For each theoretical result, does the paper provide the full set of assumptions and a complete (and correct) proof?
    \item[] Answer: \answerTODO{} % Replace by \answerYes{}, \answerNo{}, or \answerNA{}.
    \item[] Justification: \justificationTODO{}
    \item[] Guidelines:
    \begin{itemize}
        \item The answer \answerNA{} means that the paper does not include theoretical results. 
        \item All the theorems, formulas, and proofs in the paper should be numbered and cross-referenced.
        \item All assumptions should be clearly stated or referenced in the statement of any theorems.
        \item The proofs can either appear in the main paper or the supplemental material, but if they appear in the supplemental material, the authors are encouraged to provide a short proof sketch to provide intuition. 
        \item Inversely, any informal proof provided in the core of the paper should be complemented by formal proofs provided in appendix or supplemental material.
        \item Theorems and Lemmas that the proof relies upon should be properly referenced. 
    \end{itemize}

    \item {\bf Experimental result reproducibility}
    \item[] Question: Does the paper fully disclose all the information needed to reproduce the main experimental results of the paper to the extent that it affects the main claims and/or conclusions of the paper (regardless of whether the code and data are provided or not)?
    \item[] Answer: \answerTODO{} % Replace by \answerYes{}, \answerNo{}, or \answerNA{}.
    \item[] Justification: \justificationTODO{}
    \item[] Guidelines:
    \begin{itemize}
        \item The answer \answerNA{} means that the paper does not include experiments.
        \item If the paper includes experiments, a \answerNo{} answer to this question will not be perceived well by the reviewers: Making the paper reproducible is important, regardless of whether the code and data are provided or not.
        \item If the contribution is a dataset and\slash or model, the authors should describe the steps taken to make their results reproducible or verifiable. 
        \item Depending on the contribution, reproducibility can be accomplished in various ways. For example, if the contribution is a novel architecture, describing the architecture fully might suffice, or if the contribution is a specific model and empirical evaluation, it may be necessary to either make it possible for others to replicate the model with the same dataset, or provide access to the model. In general. releasing code and data is often one good way to accomplish this, but reproducibility can also be provided via detailed instructions for how to replicate the results, access to a hosted model (e.g., in the case of a large language model), releasing of a model checkpoint, or other means that are appropriate to the research performed.
        \item While NeurIPS does not require releasing code, the conference does require all submissions to provide some reasonable avenue for reproducibility, which may depend on the nature of the contribution. For example
        \begin{enumerate}
            \item If the contribution is primarily a new algorithm, the paper should make it clear how to reproduce that algorithm.
            \item If the contribution is primarily a new model architecture, the paper should describe the architecture clearly and fully.
            \item If the contribution is a new model (e.g., a large language model), then there should either be a way to access this model for reproducing the results or a way to reproduce the model (e.g., with an open-source dataset or instructions for how to construct the dataset).
            \item We recognize that reproducibility may be tricky in some cases, in which case authors are welcome to describe the particular way they provide for reproducibility. In the case of closed-source models, it may be that access to the model is limited in some way (e.g., to registered users), but it should be possible for other researchers to have some path to reproducing or verifying the results.
        \end{enumerate}
    \end{itemize}


\item {\bf Open access to data and code}
    \item[] Question: Does the paper provide open access to the data and code, with sufficient instructions to faithfully reproduce the main experimental results, as described in supplemental material?
    \item[] Answer: \answerTODO{} % Replace by \answerYes{}, \answerNo{}, or \answerNA{}.
    \item[] Justification: \justificationTODO{}
    \item[] Guidelines:
    \begin{itemize}
        \item The answer \answerNA{} means that paper does not include experiments requiring code.
        \item Please see the NeurIPS code and data submission guidelines (\url{https://neurips.cc/public/guides/CodeSubmissionPolicy}) for more details.
        \item While we encourage the release of code and data, we understand that this might not be possible, so \answerNo{} is an acceptable answer. Papers cannot be rejected simply for not including code, unless this is central to the contribution (e.g., for a new open-source benchmark).
        \item The instructions should contain the exact command and environment needed to run to reproduce the results. See the NeurIPS code and data submission guidelines (\url{https://neurips.cc/public/guides/CodeSubmissionPolicy}) for more details.
        \item The authors should provide instructions on data access and preparation, including how to access the raw data, preprocessed data, intermediate data, and generated data, etc.
        \item The authors should provide scripts to reproduce all experimental results for the new proposed method and baselines. If only a subset of experiments are reproducible, they should state which ones are omitted from the script and why.
        \item At submission time, to preserve anonymity, the authors should release anonymized versions (if applicable).
        \item Providing as much information as possible in supplemental material (appended to the paper) is recommended, but including URLs to data and code is permitted.
    \end{itemize}


\item {\bf Experimental setting/details}
    \item[] Question: Does the paper specify all the training and test details (e.g., data splits, hyperparameters, how they were chosen, type of optimizer) necessary to understand the results?
    \item[] Answer: \answerTODO{} % Replace by \answerYes{}, \answerNo{}, or \answerNA{}.
    \item[] Justification: \justificationTODO{}
    \item[] Guidelines:
    \begin{itemize}
        \item The answer \answerNA{} means that the paper does not include experiments.
        \item The experimental setting should be presented in the core of the paper to a level of detail that is necessary to appreciate the results and make sense of them.
        \item The full details can be provided either with the code, in appendix, or as supplemental material.
    \end{itemize}

\item {\bf Experiment statistical significance}
    \item[] Question: Does the paper report error bars suitably and correctly defined or other appropriate information about the statistical significance of the experiments?
    \item[] Answer: \answerTODO{} % Replace by \answerYes{}, \answerNo{}, or \answerNA{}.
    \item[] Justification: \justificationTODO{}
    \item[] Guidelines:
    \begin{itemize}
        \item The answer \answerNA{} means that the paper does not include experiments.
        \item The authors should answer \answerYes{} if the results are accompanied by error bars, confidence intervals, or statistical significance tests, at least for the experiments that support the main claims of the paper.
        \item The factors of variability that the error bars are capturing should be clearly stated (for example, train/test split, initialization, random drawing of some parameter, or overall run with given experimental conditions).
        \item The method for calculating the error bars should be explained (closed form formula, call to a library function, bootstrap, etc.)
        \item The assumptions made should be given (e.g., Normally distributed errors).
        \item It should be clear whether the error bar is the standard deviation or the standard error of the mean.
        \item It is OK to report 1-sigma error bars, but one should state it. The authors should preferably report a 2-sigma error bar than state that they have a 96\% CI, if the hypothesis of Normality of errors is not verified.
        \item For asymmetric distributions, the authors should be careful not to show in tables or figures symmetric error bars that would yield results that are out of range (e.g., negative error rates).
        \item If error bars are reported in tables or plots, the authors should explain in the text how they were calculated and reference the corresponding figures or tables in the text.
    \end{itemize}

\item {\bf Experiments compute resources}
    \item[] Question: For each experiment, does the paper provide sufficient information on the computer resources (type of compute workers, memory, time of execution) needed to reproduce the experiments?
    \item[] Answer: \answerTODO{} % Replace by \answerYes{}, \answerNo{}, or \answerNA{}.
    \item[] Justification: \justificationTODO{}
    \item[] Guidelines:
    \begin{itemize}
        \item The answer \answerNA{} means that the paper does not include experiments.
        \item The paper should indicate the type of compute workers CPU or GPU, internal cluster, or cloud provider, including relevant memory and storage.
        \item The paper should provide the amount of compute required for each of the individual experimental runs as well as estimate the total compute. 
        \item The paper should disclose whether the full research project required more compute than the experiments reported in the paper (e.g., preliminary or failed experiments that didn't make it into the paper). 
    \end{itemize}
    
\item {\bf Code of ethics}
    \item[] Question: Does the research conducted in the paper conform, in every respect, with the NeurIPS Code of Ethics \url{https://neurips.cc/public/EthicsGuidelines}?
    \item[] Answer: \answerTODO{} % Replace by \answerYes{}, \answerNo{}, or \answerNA{}.
    \item[] Justification: \justificationTODO{}
    \item[] Guidelines:
    \begin{itemize}
        \item The answer \answerNA{} means that the authors have not reviewed the NeurIPS Code of Ethics.
        \item If the authors answer \answerNo, they should explain the special circumstances that require a deviation from the Code of Ethics.
        \item The authors should make sure to preserve anonymity (e.g., if there is a special consideration due to laws or regulations in their jurisdiction).
    \end{itemize}


\item {\bf Broader impacts}
    \item[] Question: Does the paper discuss both potential positive societal impacts and negative societal impacts of the work performed?
    \item[] Answer: \answerTODO{} % Replace by \answerYes{}, \answerNo{}, or \answerNA{}.
    \item[] Justification: \justificationTODO{}
    \item[] Guidelines:
    \begin{itemize}
        \item The answer \answerNA{} means that there is no societal impact of the work performed.
        \item If the authors answer \answerNA{} or \answerNo, they should explain why their work has no societal impact or why the paper does not address societal impact.
        \item Examples of negative societal impacts include potential malicious or unintended uses (e.g., disinformation, generating fake profiles, surveillance), fairness considerations (e.g., deployment of technologies that could make decisions that unfairly impact specific groups), privacy considerations, and security considerations.
        \item The conference expects that many papers will be foundational research and not tied to particular applications, let alone deployments. However, if there is a direct path to any negative applications, the authors should point it out. For example, it is legitimate to point out that an improvement in the quality of generative models could be used to generate Deepfakes for disinformation. On the other hand, it is not needed to point out that a generic algorithm for optimizing neural networks could enable people to train models that generate Deepfakes faster.
        \item The authors should consider possible harms that could arise when the technology is being used as intended and functioning correctly, harms that could arise when the technology is being used as intended but gives incorrect results, and harms following from (intentional or unintentional) misuse of the technology.
        \item If there are negative societal impacts, the authors could also discuss possible mitigation strategies (e.g., gated release of models, providing defenses in addition to attacks, mechanisms for monitoring misuse, mechanisms to monitor how a system learns from feedback over time, improving the efficiency and accessibility of ML).
    \end{itemize}
    
\item {\bf Safeguards}
    \item[] Question: Does the paper describe safeguards that have been put in place for responsible release of data or models that have a high risk for misuse (e.g., pre-trained language models, image generators, or scraped datasets)?
    \item[] Answer: \answerTODO{} % Replace by \answerYes{}, \answerNo{}, or \answerNA{}.
    \item[] Justification: \justificationTODO{}
    \item[] Guidelines:
    \begin{itemize}
        \item The answer \answerNA{} means that the paper poses no such risks.
        \item Released models that have a high risk for misuse or dual-use should be released with necessary safeguards to allow for controlled use of the model, for example by requiring that users adhere to usage guidelines or restrictions to access the model or implementing safety filters. 
        \item Datasets that have been scraped from the Internet could pose safety risks. The authors should describe how they avoided releasing unsafe images.
        \item We recognize that providing effective safeguards is challenging, and many papers do not require this, but we encourage authors to take this into account and make a best faith effort.
    \end{itemize}

\item {\bf Licenses for existing assets}
    \item[] Question: Are the creators or original owners of assets (e.g., code, data, models), used in the paper, properly credited and are the license and terms of use explicitly mentioned and properly respected?
    \item[] Answer: \answerTODO{} % Replace by \answerYes{}, \answerNo{}, or \answerNA{}.
    \item[] Justification: \justificationTODO{}
    \item[] Guidelines:
    \begin{itemize}
        \item The answer \answerNA{} means that the paper does not use existing assets.
        \item The authors should cite the original paper that produced the code package or dataset.
        \item The authors should state which version of the asset is used and, if possible, include a URL.
        \item The name of the license (e.g., CC-BY 4.0) should be included for each asset.
        \item For scraped data from a particular source (e.g., website), the copyright and terms of service of that source should be provided.
        \item If assets are released, the license, copyright information, and terms of use in the package should be provided. For popular datasets, \url{paperswithcode.com/datasets} has curated licenses for some datasets. Their licensing guide can help determine the license of a dataset.
        \item For existing datasets that are re-packaged, both the original license and the license of the derived asset (if it has changed) should be provided.
        \item If this information is not available online, the authors are encouraged to reach out to the asset's creators.
    \end{itemize}

\item {\bf New assets}
    \item[] Question: Are new assets introduced in the paper well documented and is the documentation provided alongside the assets?
    \item[] Answer: \answerTODO{} % Replace by \answerYes{}, \answerNo{}, or \answerNA{}.
    \item[] Justification: \justificationTODO{}
    \item[] Guidelines:
    \begin{itemize}
        \item The answer \answerNA{} means that the paper does not release new assets.
        \item Researchers should communicate the details of the dataset\slash code\slash model as part of their submissions via structured templates. This includes details about training, license, limitations, etc. 
        \item The paper should discuss whether and how consent was obtained from people whose asset is used.
        \item At submission time, remember to anonymize your assets (if applicable). You can either create an anonymized URL or include an anonymized zip file.
    \end{itemize}

\item {\bf Crowdsourcing and research with human subjects}
    \item[] Question: For crowdsourcing experiments and research with human subjects, does the paper include the full text of instructions given to participants and screenshots, if applicable, as well as details about compensation (if any)? 
    \item[] Answer: \answerTODO{} % Replace by \answerYes{}, \answerNo{}, or \answerNA{}.
    \item[] Justification: \justificationTODO{}
    \item[] Guidelines:
    \begin{itemize}
        \item The answer \answerNA{} means that the paper does not involve crowdsourcing nor research with human subjects.
        \item Including this information in the supplemental material is fine, but if the main contribution of the paper involves human subjects, then as much detail as possible should be included in the main paper. 
        \item According to the NeurIPS Code of Ethics, workers involved in data collection, curation, or other labor should be paid at least the minimum wage in the country of the data collector. 
    \end{itemize}

\item {\bf Institutional review board (IRB) approvals or equivalent for research with human subjects}
    \item[] Question: Does the paper describe potential risks incurred by study participants, whether such risks were disclosed to the subjects, and whether Institutional Review Board (IRB) approvals (or an equivalent approval/review based on the requirements of your country or institution) were obtained?
    \item[] Answer: \answerTODO{} % Replace by \answerYes{}, \answerNo{}, or \answerNA{}.
    \item[] Justification: \justificationTODO{}
    \item[] Guidelines:
    \begin{itemize}
        \item The answer \answerNA{} means that the paper does not involve crowdsourcing nor research with human subjects.
        \item Depending on the country in which research is conducted, IRB approval (or equivalent) may be required for any human subjects research. If you obtained IRB approval, you should clearly state this in the paper. 
        \item We recognize that the procedures for this may vary significantly between institutions and locations, and we expect authors to adhere to the NeurIPS Code of Ethics and the guidelines for their institution. 
        \item For initial submissions, do not include any information that would break anonymity (if applicable), such as the institution conducting the review.
    \end{itemize}

\item {\bf Declaration of LLM usage}
    \item[] Question: Does the paper describe the usage of LLMs if it is an important, original, or non-standard component of the core methods in this research? Note that if the LLM is used only for writing, editing, or formatting purposes and does \emph{not} impact the core methodology, scientific rigor, or originality of the research, declaration is not required.
    %this research? 
    \item[] Answer: \answerTODO{} % Replace by \answerYes{}, \answerNo{}, or \answerNA{}.
    \item[] Justification: \justificationTODO{}
    \item[] Guidelines:
    \begin{itemize}
        \item The answer \answerNA{} means that the core method development in this research does not involve LLMs as any important, original, or non-standard components.
        \item Please refer to our LLM policy in the NeurIPS handbook for what should or should not be described.
    \end{itemize}

\end{enumerate}

\end{document}
 
